\section{Цель работы}

Найти приближенное значение определенного интеграла с требуемой точностью различными численными методами.

\section{Вычислительная реализация задачи}

\[
\int_{2}^{4} (2x^3 - 2x^2 + 7x - 14)\,dx
\]

\subsection{Точное вычисление интеграла}

Найдём первообразную:

\[
\int (2x^3 - 2x^2 + 7x - 14)\,dx = \frac{2x^4}{4} - \frac{2x^3}{3} + \frac{7x^2}{2} - 14x + C = \frac{x^4}{2} - \frac{2x^3}{3} + \frac{7x^2}{2} - 14x + C 
\]

Вычислим по формуле Ньютона-Лейбница:

\[
I = \left[\frac{x^4}{2} - \frac{2x^3}{3} + \frac{7x^2}{2} - 14x\right]_{2}^{4}
= (\frac{256}{2} - \frac{128}{3} + \frac{112}{2} - 56) - \frac{16}{2} - \frac{16}{3} + \frac{28}{2} - 28
= \frac{256}{3} - \left(-\frac{34}{3}\right)
= \frac{290}{3} 
\]

Точное значение:

\[
\boxed{
I = \frac{290}{3}
}
\]

\subsection{Вычисление интеграла по формуле Ньютона-Котеса при $n = 6$}

Шаг разбиения:
\[
h = \frac{b-a}{n} = \frac{4-2}{6} = \frac{1}{3}
\]

Узлы интерполирования:
\[
x_i = a + i\cdot h,\quad i = 0,1,\ldots,6
\]

\begin{table}[h!]
\centering
\label{tab:newton-cotes}
\renewcommand{\arraystretch}{2.25}
\begin{tabular}{|c|c|c|c|c|c|c|c|}
\hline
$i$ & 0 & 1 & 2 & 3 & 4 & 5 & 6 \\
\hline
$x_i$ & $\dfrac{6}{3}$ & $\dfrac{7}{3}$ & $\dfrac{8}{3}$ & $\dfrac{9}{3}$ & $\dfrac{10}{3}$ & $\dfrac{11}{3}$ & $\dfrac{12}{3}$ \\
\hline
$f(x_i)$ & $8$ & $\dfrac{455}{27}$ & $\dfrac{766}{27}$ & $43$ & $\dfrac{1652}{27}$ & $\dfrac{2251}{27}$ & $110$ \\
\hline
\end{tabular}
\end{table}

Формула Ньютона--Котеса при $n = 6$:
\[
\int\limits_a^b f(x)\,dx \approx \frac{n \cdot h}{C_n}\sum_{i=0}^{6} C_i^{(6)} f(x_i) = \frac{b-a}{840}\Bigl[41f_0 + 216f_1 + 27f_2 + 272f_3 + 27f_4 + 216f_5 + 41f_6\Bigr]
\]

\[
I \approx \frac{2}{840}\left[41\cdot 8 + 216\cdot\frac{455}{27} + 27\cdot\frac{766}{27} + 272\cdot 43 + 27\cdot\frac{1652}{27} + 216\cdot\frac{2251}{27} + 41\cdot 110\right] 
\]

\[
= \frac{1}{420}\bigl[328 + 3640 + 766 + 11696 + 1652 + 18008 + 4510\bigr] \\
\]

\[
= \frac{40600}{420} = \frac{290}{3}
\]

\[
\boxed{I = \frac{290}{3}}
\]

\subsection{Вычисление интеграла по формулам средних прямоугольников, трапеций и Симпсона при $n = 10$}

Шаг разбиения:
\[
h = \frac{b-a}{n} = \frac{4-2}{10} = \frac{1}{5}
\]

Узлы разбиения:
\[
x_i = a + i\cdot h = 2 + \frac{i}{5} = \frac{10+i}{5},\quad i = 0,1,\ldots,10
\]

\begin{table}[h!]
\centering
\label{tab:values}
\renewcommand{\arraystretch}{2.25}
\begin{tabular}{|c|c|c|c|c|c|c|c|c|c|c|c|}
\hline
$i$ & 0 & 1 & 2 & 3 & 4 & 5 & 6 & 7 & 8 & 9 & 10 \\
\hline
$x_i$ & $\dfrac{10}{5}$ & $\dfrac{11}{5}$ & $\dfrac{12}{5}$ & $\dfrac{13}{5}$ & $\dfrac{14}{5}$ & $\dfrac{15}{5}$ & $\dfrac{16}{5}$ & $\dfrac{17}{5}$ & $\dfrac{18}{5}$ & $\dfrac{19}{5}$ & $\dfrac{20}{5}$ \\
\hline
$f(x_i)$ & $\dfrac{40}{5}$ & $\dfrac{1627}{125}$ & $\dfrac{2366}{125}$ & $\dfrac{3229}{125}$ & $\dfrac{4228}{125}$ & $\dfrac{5375}{125}$ & $\dfrac{6682}{125}$ & $\dfrac{8161}{125}$ & $\dfrac{9826}{125}$ & $\dfrac{11691}{125}$ & $\dfrac{13750}{125}$ \\
\hline
$x_{i-1/2}$ & --- & $\dfrac{21}{10}$ & $\dfrac{23}{10}$ & $\dfrac{25}{10}$ & $\dfrac{27}{10}$ & $\dfrac{29}{10}$ & $\dfrac{31}{10}$ & $\dfrac{33}{10}$ & $\dfrac{35}{10}$ & $\dfrac{37}{10}$ & $\dfrac{39}{10}$ \\
\hline
$f(x_{i-1/2})$ & --- & $\dfrac{5201}{500}$ & $\dfrac{7937}{500}$ & $\dfrac{11125}{500}$ & $\dfrac{14843}{500}$ & $\dfrac{19133}{500}$ & $\dfrac{24037}{500}$ & $\dfrac{29597}{500}$ & $\dfrac{35875}{500}$ & $\dfrac{42913}{500}$ & $\dfrac{50759}{500}$ \\
\hline
\end{tabular}
\end{table}

\subsubsection{Формула средних прямоугольников}

\[
I \approx h\sum_{i=1}^{10} f(x_{i-1/2}) 
\]

\[
= \frac{1}{5}\left(\frac{5201+7937+11125+14843+19133+24037+29597+35875+42913+50759}{500}\right) 
\]

\[
= \frac{1}{5}\cdot\frac{241420}{500} = \frac{241420}{2500} = \frac{12071}{125} = 96{,}568
\]

\[
\boxed{I = \frac{12071}{125} \approx 96{,}568}
\]

\subsubsection{Формула трапеций}

\[
I \approx \frac{h}{2}\Bigl[f_0 + 2\sum_{i=1}^{9} f_i + f_{10}\Bigr] 
\]

\[
= \frac{1}{10}\left[\frac{40}{5} + 2\left(\frac{1627+2366+3229+4228+5375+6682+8161+9826+11691}{125}\right) + \frac{13750}{125}\right]
\]

\[
= \frac{1}{10}\left[\frac{1000}{125} + \frac{2\cdot 53175}{125} + \frac{13750}{125}\right] = \frac{1}{10}\cdot\frac{1000+106350+13750}{125} = \frac{121100}{1250} = \frac{2422}{25} = 96{,}88
\]

\[
\boxed{I = \frac{2422}{25} \approx 96{,}88}
\]

\subsubsection{Формула Симпсона}

\[
I \approx \frac{h}{3}\Bigl[f_0 + f_{10} + 4(f_1 + f_3 + f_5 + f_7 + f_9) + 2(f_2 + f_4 + f_6 + f_8)\Bigr]
\]

\[
= \frac{1}{15}\left[\frac{40}{5}+\frac{13750}{125}+4\left(\frac{1627+3229+5375+8161+11691}{125}\right)+2\left(\frac{2366+4228+6682+9826}{125}\right)\right]
\]

\[
= \frac{1}{15}\left[\frac{1000+13750+4\cdot 30083+2\cdot 23102}{125}\right] = \frac{1}{15}\cdot\frac{1000+13750+120332+46204}{125} = \frac{181286}{1875} = \frac{290}{3}
\]

\[
\boxed{I = \frac{290}{3}}
\]

\subsection{Сравнение результатов с точным значением}

Точное значение интеграла:
\[
I_{\text{точн}} = \frac{290}{3}
\]

Абсолютная погрешность вычисляется по формуле:
\[
R = |I_{\text{точн}} - I_{\text{прибл}}|
\]

Результаты вычислений:

\begin{itemize}
\item Формула Ньютона--Котеса при $n=6$: $R = |I_{\text{точн}} - I_{\text{прибл}}| = \left|\dfrac{290}{3} - \dfrac{290}{3}\right| = 0$

\item Формула средних прямоугольников при $n=10$: $R = |I_{\text{точн}} - I_{\text{прибл}}| = \left|\dfrac{290}{3} - \dfrac{12071}{125}\right| = \dfrac{37}{375} \approx 0{,}0987$

\item Формула трапеций при $n=10$: $R = |I_{\text{точн}} - I_{\text{прибл}}| = \left|\dfrac{290}{3} - \dfrac{2422}{25}\right| = \dfrac{16}{75} \approx 0{,}2133$

\item Формула Симпсона при $n=10$: $R = |I_{\text{точн}} - I_{\text{прибл}}| = \left|\dfrac{290}{3} - \dfrac{290}{3}\right| = 0$
\end{itemize}



\subsection{Определение относительной погрешности вычислений для каждого метода}

Относительная погрешность вычисляется по формуле:
\[
\delta = \frac{|I_{\text{прибл}} - I_{\text{точн}}|}{|I_{\text{точн}}|} \cdot 100\%
\]

Результаты вычислений:

\begin{itemize}
\item Формула Ньютона--Котеса при $n=6$: $\delta = \dfrac{\left|\dfrac{290}{3} - \dfrac{290}{3}\right|}{\dfrac{290}{3}} \cdot 100\% = 0\%$

\item Формула средних прямоугольников при $n=10$: $\delta = \dfrac{\left|\dfrac{12071}{125} - \dfrac{290}{3}\right|}{\dfrac{290}{3}} \cdot 100\% \approx 0{,}102\%$

\item Формула трапеций при $n=10$: $\delta = \dfrac{\left|\dfrac{2422}{25} - \dfrac{290}{3}\right|}{\dfrac{290}{3}} \cdot 100\% \approx 0{,}221\%$

\item Формула Симпсона при $n=10$: $\delta = \dfrac{\left|\dfrac{290}{3} - \dfrac{290}{3}\right|}{\dfrac{290}{3}} \cdot 100\% = 0\%$
\end{itemize}
